\documentclass[a4paper,10pt]{article}
\usepackage{xeCJK}
\setCJKmainfont{Hiragino Mincho ProN}
\setCJKsansfont{Hiragino Sans}
\setCJKmonofont{Hiragino Sans}

\usepackage[top=12.5mm,bottom=12.5mm,left=13mm,right=13mm]{geometry}
\usepackage{amsmath,amssymb,bm}
\usepackage{multicol}
\usepackage{enumitem}
\usepackage{titlesec}
\usepackage[hidelinks]{hyperref}

\pagestyle{empty}
\setlength{\parindent}{1em}
\setlength{\parskip}{0pt}
\setlength{\columnsep}{5.8mm}
\renewcommand{\baselinestretch}{0.95}

\titleformat{\section}
  {\normalfont\bfseries\normalsize}
  {\thesection.}{0.45em}{}
\titlespacing*{\section}{0pt}{0.45\baselineskip}{0.12\baselineskip}

\newcommand{\E}{\mathbb E}
\newcommand{\Prb}{\mathbb P}
\newcommand{\indep}{\mathrel{\perp\!\!\!\perp}}
\newcommand{\Y}{\bm Y}
\newcommand{\Z}{\bm Z}

\newcommand{\papertitle}{欠測IVと潜在変数モデリングによる非無作為欠測の識別}
\newcommand{\subtitle}{構造的MNAR・supported blocks・潜在測定核の1ページ要約}
\newcommand{\authorblock}{%
  \begin{tabular}{c}
    本多 将大\(^{1,2}\)、星野 崇宏\(^{1,2}\)、大津 泰介\(^{1,3}\)\\[-0.1mm]
    {\scriptsize 1. 慶應義塾大学、2. 理化学研究所、3. London School of Economics and Political Science}
  \end{tabular}
}

\begin{document}

\vspace*{-9mm}
\begin{center}
  {\large\bfseries \papertitle\par}
  \vspace{0.2mm}
  {\small \subtitle\par}
  \vspace{0.5mm}
  {\footnotesize \authorblock}
\end{center}
\vspace{-3.5mm}

\footnotesize
\setlength{\baselineskip}{9.7pt}
\setlength{\abovedisplayskip}{1.2pt}
\setlength{\belowdisplayskip}{1.2pt}
\setlength{\abovedisplayshortskip}{0.8pt}
\setlength{\belowdisplayshortskip}{0.8pt}
\setlength{\jot}{0pt}
\begin{multicols}{2}

\section{何を解きたいか}
個体 \(i\) の添字を省略し，潜在的測定値と観測指標を
\[
  Y=(Y_1,\ldots,Y_m),\qquad
  D=(D_1,\ldots,D_m)\in\{0,1\}^m
\]
とする。\(Y_j\) は \(D_j=1\) のときだけ観測される。特に exact top-\(K\) 型では
\[
  \sum_{j=1}^m D_j=K<m
\]
であり，完全ケース \(D=\bm1_m\) は構造的に起こらない。したがって通常の「全項目が観測された人だけで潜在因子を推定する」戦略は出発点から使えない。

潜在測定モデルは
\[
\begin{gathered}
F\mid(W,X)\sim Q_F,\qquad
Y_j\mid(F,X)\sim Q_j,\\
Y_1,\ldots,Y_m\indep\mid(F,X).
\end{gathered}
\]
で表す。ここで \(W\) は潜在分布 \(P(F\mid W,X)\) を動かす常時観測変数であり，\(Z_S\) は block \(S\) の欠測を補正する shadow IV である。重要なのは，\(W\) と \(Z_S\) の役割を分ける点である。仮に
\[
  D\indep W\mid(Y,F,X)
\]
でも，\(F\) を積分すると \(P(F\mid Y,W,X)\) が \(W\) に依存し得るため，一般に \(D\indep W\mid(Y,X)\) は従わない。

\section{基本アイデア}
noteの主張は，構造的に完全ケースがない多変量MNARを，二段階で識別するというものである。

\[
\begin{aligned}
O=(X,W,Z,D,Y_D)&\Longrightarrow P(Y_S,Z_S,C_S),\\
&\Longrightarrow \{P(F\mid W,X),P(Y_j\mid F,X)\}.
\end{aligned}
\]

\noindent\textbf{第1段階：supported blockの回復.}
集合 \(S\subset\{1,\ldots,m\}\) について
\[
  R_S=\prod_{j\in S}D_j,\qquad Y_S=(Y_j:j\in S)
\]
とおく。\(R_S=1\) なら \(Y_S\) は同時観測される。文脈 \(C_S\) の下で
\[
\begin{gathered}
\pi_S(y,c)=\Prb(R_S=1\mid Y_S=y,C_S=c)>0,\\
R_S\indep Z_S\mid(Y_S,C_S)
\end{gathered}
\]
を仮定し，さらに \(Y_S\) が \(Z_S\) に対して十分豊かに動くという完備性を置く。このとき d'Haultfoeuille 型の条件付きモーメント
\[
  \E\left\{\frac{R_S}{q(Y_S,C_S)}\mid Z_S,C_S\right\}=1
\]
の唯一の正値解が \(q=\pi_S\) となる。したがって
\[
  \E\{h(Y_S,Z_S,C_S)\}
  =
  \E\left\{
    \frac{R_S}{\pi_S(Y_S,C_S)}h(Y_S,Z_S,C_S)
  \right\}
\]
により，観測された block から \(P(Y_S,Z_S,C_S)\) を回復できる。分布全体が不要な場合は，特定の平均・セル確率だけを representer bridge で回復すればよい。

\columnbreak

\section{潜在因子への接続}
\noindent\textbf{第2段階：blockを潜在測定モデルでつなぐ.}
第1段階で回復した複数の低次元分布を，潜在因子 \(F\) を共有する測定モデルとして接続する。
有限潜在クラス \(F\in\{1,\ldots,r\}\) では
\[
\begin{aligned}
M_j(x)&=\bigl[\Prb(Y_j=y\mid F=f,X=x)\bigr]_{y,f},\\
\lambda_f(w,x)&=\Prb(F=f\mid W=w,X=x).
\end{aligned}
\]
とおく。supported な anchor pair または anchor triple と，\(W\) による混合比率の変動が十分に rank を持てば，回復済みテンソルを分解して
\[
  \{M_j(x):j=1,\ldots,m\},\qquad
  \{\lambda_f(w,x):f=1,\ldots,r\}
\]
を識別できる。結果として，一度も同時観測されない \(|A|>K\) の項目集合についても
\[
  \Prb(Y_A=y_A\mid W=w,X=x)
  =
  \sum_{f=1}^r \lambda_f(w,x)
  \prod_{j\in A} M_j(y_j,f;x)
\]
を外挿できる。

連続 \(F\) では，行列rankの代わりに条件付き期待値作用素の injectivity と同時対角化を用いる。ただし平均・分散・符号だけでは非線形な再パラメータ化を完全には除けないため，主張は追加正規化と作用素の定義域条件の下に限定する。

\section{JoE 2010との違い}
d'Haultfoeuille (2010) は，スカラー欠測アウトカムで shadow IV と完備性により選択確率と観測変数の同時分布を識別する。本noteが狙う拡張は，単にその結果を繰り返すことではない。

\begin{itemize}[leftmargin=1.4em,itemsep=0pt,topsep=0pt]
\item 完全ケースが構造的に存在しない top-\(K\) / pattern MNAR を扱う。
\item supported block ごとに欠測IVで law を回復する。
\item 回復した block laws をテンソル分解・作用素分解で接続し，潜在分布と測定核を識別する。
\item \(K=1\)，anchor support 欠如，rank deficiency，shadow violation は識別境界として明示する。
\end{itemize}

したがって防御可能な新規性は，「潜在変数MNAR一般を解く」ことではなく，\textbf{blockwise shadow recovery と latent measurement identification を exact top-\(K\) support geometry 上で統合する}点にある。

\section{実行前チェック}
応用に進む前に，少なくとも次を監査する。
\begin{enumerate}[label=\arabic*.,leftmargin=1.6em,itemsep=0pt,topsep=0pt]
\item \(Z_S\) は全員について，選択より前に測定されているか。
\item \(Y_S,C_S\) を固定した後，\(Z_S\) が \(R_S\) に直接影響しない説明が可能か。
\item 必要な supported pair/triple の同時観測確率が十分か。
\item weighted tensor/operator の最小特異値がゼロから離れているか。
\item latent class の数，scale，sign，label の正規化を明示したか。
\item protocol missing-by-design，死亡 truncation，administrative censoring と MNAR を区別したか。
\end{enumerate}

\vspace{0.3mm}
\noindent\textbf{まとめ.}
noteの中心は，\(F\) を単に積分して欠測IV条件を得るのではなく，supported block を先に回復し，その後で潜在測定モデルとして block を接続する点にある。この二段階を推定・感度分析・医学応用と一体化できれば，JoE 2010 の単純な多変量化を超える研究になる。

\vspace{0.2mm}
\noindent{\scriptsize 主要文献：d'Haultfoeuille (2010, JoE), Hu and Schennach (2008, Econometrica), Li et al. (2023, JRSS-B).}

\end{multicols}
\end{document}
